% GENERATED FILE. Edit canonical lessons, then run npm run generate:books.
% canonical-source-hash: fnv1a64:e6e0db17febbd4f0
% canonical-lessons: ES-C729-terco, ES-C729-perezoso, ES-C729-pereza, ES-C729-entrometido, ES-C729-curiosidad

\chapter{Defectos de cada día}
\label{ch:es-defectos-de-cada-dia}
\hlchaptermodality{729}

\begin{chapteropening}
\textbf{By the end of this chapter:} \emph{I can talk about stubborn, lazy, nosy and curious people.}

Complete the last lesson of chapter 729: I can talk about stubborn, lazy, nosy and curious people.
\end{chapteropening}

\section[terco]{terco — stubborn}

\label{lesson:ES-C729-terco}



\begin{quote}
Before the new one: say the Spanish for skilful, then the Spanish for wit.
\end{quote}

\subsection*{You'll want to know: terco}
\textbf{terco} — \textquotedblleft{}stubborn\textquotedblright{}.

\begin{quote}
Es tan terco como su padre.
\end{quote}

\begin{quote}
\emph{He is as stubborn as his father.}
\end{quote}

\emph{Cabezota} is the colloquial word in Spain.

\begin{cousinweb}
For \emph{terco} the origin is not settled. It is found in old Spanish, and no Latin or Germanic source has been proved.
\end{cousinweb}

\begin{grammarlens}[title={What you've built}]
One more word for telling, explaining and describing at B1.
\end{grammarlens}

\subsection*{Your turn}
\begin{itemize}
\raggedright
  \item \emph{Say it:} \emph{terco}
  \item \emph{Say it:} \emph{terco}, once more
  \item \emph{Say it:} a sentence of your own with \emph{terco}
  \item \emph{Recall:} say the Spanish for skilful, then the Spanish for wit, then say \emph{terco} again
\end{itemize}

\begin{tcolorbox}[breakable,colback=teal!4,colframe=teal!35!black,title={Before you move on}]
What does terco mean? (\textquotedblleft{}Stubborn\textquotedblright{}.) Say it once more.
\end{tcolorbox}
\section[perezoso]{perezoso — lazy}

\label{lesson:ES-C729-perezoso}



\begin{quote}
Before the new one: say the Spanish for wit, then the Spanish for stubborn.
\end{quote}

\subsection*{You'll want to know: perezoso}
\textbf{perezoso} — \textquotedblleft{}lazy\textquotedblright{}.

\begin{quote}
Los domingos soy muy perezoso.
\end{quote}

\begin{quote}
\emph{On Sundays I am very lazy.}
\end{quote}

\emph{Un perezoso} is also the animal, a sloth.

\begin{cousinweb}
\emph{perezoso} is built on \emph{pereza} \textquotedblleft{}laziness\textquotedblright{} with the ending \emph{-oso}; \emph{pereza} is usually traced to Latin \emph{pigritia}.
\end{cousinweb}

\begin{grammarlens}[title={What you've built}]
One more word for telling, explaining and describing at B1.
\end{grammarlens}

\subsection*{Your turn}
\begin{itemize}
\raggedright
  \item \emph{Say it:} \emph{perezoso}
  \item \emph{Say it:} \emph{perezoso}, once more
  \item \emph{Say it:} a sentence of your own with \emph{perezoso}
  \item \emph{Recall:} say the Spanish for wit, then the Spanish for stubborn, then say \emph{perezoso} again
\end{itemize}

\begin{tcolorbox}[breakable,colback=teal!4,colframe=teal!35!black,title={Before you move on}]
What does perezoso mean? (\textquotedblleft{}Lazy\textquotedblright{}.) Say it once more.
\end{tcolorbox}
\section[la pereza]{la pereza — laziness}

\label{lesson:ES-C729-pereza}



\begin{quote}
Before the new one: say the Spanish for stubborn, then the Spanish for lazy.
\end{quote}

\subsection*{You'll want to know: la pereza}
\textbf{la pereza} — \textquotedblleft{}laziness\textquotedblright{}.

\begin{quote}
Me da pereza salir con este frío.
\end{quote}

\begin{quote}
\emph{I cannot be bothered to go out in this cold.}
\end{quote}

\emph{Me da pereza} + infinitive — I cannot be bothered to.

\begin{cousinweb}
\emph{pereza} is usually traced to Latin \emph{pigritia} \textquotedblleft{}laziness\textquotedblright{}, built on \emph{piger} \textquotedblleft{}slow, lazy\textquotedblright{}; Latin \emph{-itia} became Spanish \emph{-eza}, as in \emph{tristeza}.
\end{cousinweb}

\begin{grammarlens}[title={What you've built}]
One more word for telling, explaining and describing at B1.
\end{grammarlens}

\subsection*{Your turn}
\begin{itemize}
\raggedright
  \item \emph{Say it:} \emph{la pereza}
  \item \emph{Say it:} \emph{la pereza}, once more
  \item \emph{Say it:} a sentence of your own with \emph{la pereza}
  \item \emph{Recall:} say the Spanish for stubborn, then the Spanish for lazy, then say \emph{la pereza} again
\end{itemize}

\begin{tcolorbox}[breakable,colback=teal!4,colframe=teal!35!black,title={Before you move on}]
What does la pereza mean? (\textquotedblleft{}Laziness\textquotedblright{}.) Say it once more.
\end{tcolorbox}
\section[entrometido]{entrometido — nosy, meddling}

\label{lesson:ES-C729-entrometido}



\begin{quote}
Before the new one: say the Spanish for lazy, then the Spanish for laziness.
\end{quote}

\subsection*{You'll want to know: entrometido}
\textbf{entrometido} — \textquotedblleft{}nosy, meddling\textquotedblright{}.

\begin{quote}
Mi vecina es muy entrometida.
\end{quote}

\begin{quote}
\emph{My neighbour is very nosy.}
\end{quote}

\emph{Entrometerse} — to meddle. Feminine \emph{entrometida}.

\begin{cousinweb}
\emph{entrometido} is the participle of \emph{entrometer}, built on \emph{entre} \textquotedblleft{}between\textquotedblright{} and \emph{meter} \textquotedblleft{}to put\textquotedblright{}, which is inherited from Latin \emph{mittere} \textquotedblleft{}to send\textquotedblright{} (EtymDB).
\end{cousinweb}

\begin{grammarlens}[title={What you've built}]
One more word for telling, explaining and describing at B1.
\end{grammarlens}

\subsection*{Your turn}
\begin{itemize}
\raggedright
  \item \emph{Say it:} \emph{entrometido}
  \item \emph{Say it:} \emph{entrometido}, once more
  \item \emph{Say it:} a sentence of your own with \emph{entrometido}
  \item \emph{Recall:} say the Spanish for lazy, then the Spanish for laziness, then say \emph{entrometido} again
\end{itemize}

\begin{tcolorbox}[breakable,colback=teal!4,colframe=teal!35!black,title={Before you move on}]
What does entrometido mean? (\textquotedblleft{}Nosy, meddling\textquotedblright{}.) Say it once more.
\end{tcolorbox}
\section[la curiosidad]{la curiosidad — curiosity}

\label{lesson:ES-C729-curiosidad}



\begin{quote}
Before the new one: say the Spanish for laziness, then the Spanish for nosy.
\end{quote}

\subsection*{You'll want to know: la curiosidad}
\textbf{la curiosidad} — \textquotedblleft{}curiosity\textquotedblright{}.

\begin{quote}
Lo pregunto por curiosidad.
\end{quote}

\begin{quote}
\emph{I am asking out of curiosity.}
\end{quote}

\emph{Por curiosidad} — out of curiosity.

\begin{cousinweb}
\emph{curiosidad} was borrowed from Latin \emph{cūriōsitās} (EtymDB), built on \emph{cūriōsus}, from \emph{cūra} \textquotedblleft{}care\textquotedblright{}.
\end{cousinweb}

\begin{grammarlens}[title={What you've built}]
One more word for telling, explaining and describing at B1.
\end{grammarlens}

\subsection*{Your turn}
\begin{itemize}
\raggedright
  \item \emph{Say it:} \emph{la curiosidad}
  \item \emph{Say it:} \emph{la curiosidad}, once more
  \item \emph{Say it:} a sentence of your own with \emph{la curiosidad}
  \item \emph{Recall:} say the Spanish for laziness, then the Spanish for nosy, then say \emph{la curiosidad} again
\end{itemize}

\begin{tcolorbox}[breakable,colback=teal!4,colframe=teal!35!black,title={Before you move on}]
What does la curiosidad mean? (\textquotedblleft{}Curiosity\textquotedblright{}.) Say it once more.
\end{tcolorbox}
